% § 1 — Introduction
\section{Introduction}
\label{sec:intro}
Directive 93/13/EEC on unfair terms in consumer contracts~\cite{eu1993directive} defines unfairness by a general test---a significant imbalance in the parties' rights and obligations, contrary to good faith---and complements it with an Annex: an indicative and non-exhaustive list of seventeen kinds of terms, \gl{a} to \gl{q}, that \emph{may} be regarded as unfair. Practitioners call it the grey list. It is the closest thing European consumer law offers to an operational definition of the phenomenon, and it is written in terms of \emph{who} may do \emph{what} to \emph{whom} under \emph{which} conditions: a seller who may terminate a contract of indeterminate duration without reasonable notice \gi{g}, alter the terms unilaterally without a valid reason specified in the contract \gi{j}, or exclude liability for personal injury \gi{a}.

Automated detectors of unfair terms do not use it. Since \corpus{}~\cite{lippi2019claudette}, the reference benchmark for online terms of service (ToS), the task has been framed as sentence classification: a learned model assigns one of eight unfairness categories to each sentence, and its decision cites no legal ground. Rationale-based extensions~\cite{ruggeri2022memory,liepina2020claudettetool} attach explanations to predictions, but the explanations are learned too. Large language models prompted on the same task are neither more accurate nor more accountable~\cite{panarelli2025worth,frasheri2024llm}.

This paper asks a simple question: \textbf{what happens if the grey list is executed literally?} We turn each item into a query over a structured representation of the contract: a \emph{clause template} recording who acts, what is obliged, permitted or forbidden, and under which conditions, deadlines and remedies, attached to clauses identified by a human-validated thematic layer over the 50 \corpus{} contracts. That layer is the resource built and validated in our companion short paper~\cite{elazhar2026thematic}, which asks how reliably legal themes can be annotated; the present paper is its sequel and asks what the themes make possible, namely the execution of the law itself. Queries are written before any template is extracted and frozen before any computation on held-out contracts, so that the result is a measurement of the grey list and not of our ability to fit it. Three questions structure the paper, each with a hypothesis recorded in advance (Table~\ref{tab:rq}). \textbf{RQ-A} (Section~\ref{sec:results}): which items can be expressed, and do the frozen queries retrieve the reference labels better than the clause theme alone? \textbf{RQ-B} (Section~\ref{sec:cost}): what does this interpretability cost against learned text-only models? \textbf{RQ-C} (Section~\ref{sec:audit}): what do the false positives reveal about the benchmark?

Our contributions are: (i) an executable, frozen encoding of the grey list, with an analysis of which items can be expressed as template conditions; (ii) a pre-registered, per-item evaluation against reference labels on held-out contracts, including the gain over the clause theme alone; (iii) a cost accounting of interpretability against text-only models trained on the same split; (iv) an error analysis in which every flag is traceable to the fields and sentences that triggered it; and (v) the release of queries, hashes, templates and code.
